% !TeX root = ../papers/nc-p3.tex
\section{有限点グラフに加法群が作用しないこと}
\label{sec:finite-unipotent-exclusion}
\begin{lemma}\label{lem:maximal-tangent-no-additive}
三次元接点を持つ標準正則有限窓の射影安定群は加法群を含まない。
\end{lemma}
\begin{proof}
定理\ref{thm:maximal-tangent-complete}の一意な三次元接点の各位置を $p_i$ とする。
加法群はこの点列を保ち、$H_i=\ker p_i$ も保つ。
定理\ref{thm:maximal-tangent-ore}の内在的な三文字関係 $S_i=R_i\cap(H_i\otimes H_{i+1})$ と
その点対応の射影像も保存される。定理\ref{thm:reduced-ternary-graphs}により
各 $\mathbb P(H_i^*)$ でこの像は滑らかな三次曲線または節点だけを特異点とする被約三次曲線である。
どの許容曲線にも非自明な射影加法群は作用しない。
滑らかな種数1曲線には射影直線からの非定数射がなく、加法群の軌道写像も固有性で射影直線へ延びる。
節点有理三次曲線では正規化上の二点が固定され、横断する直線と二次曲線、
非共点三直線でも各成分上で少なくとも二つの交点が固定される。
射影直線の二点を固定する乗法群には加法群からの非自明な準同型がない。
これらの平面曲線は射影標構を含むため曲線上の恒等作用は平面上でも恒等である。

したがって冪零生成元は $H_i$ を零にし、残る文字 $Z_i$ をある $n_i\in H_i$ に送る。
同定理による各位置の $H_i$ の移送を行い、混合関係を
$Z_i h-h Z_{i+1}-d_i(h)$ と書く。微分すると
$n_i\otimes h-h\otimes n_{i+1}\in S_i$ が全 $h$ について必要である。
三文字の点グラフの $(p,\sigma p)$ で評価すれば
$n_i(p)\sigma p=n_{i+1}(\sigma p)p$ となる（代表の倍率は両辺で同じ）。
補題\ref{lem:reduced-ternary-fourth-shift}により全既約成分の一般点で $p$ と $\sigma p$ は独立である。
ゆえに $n_i$ は被約三次曲線全体で消える一次式であり零である。これで全位置の作用が消える。
\end{proof}

\begin{theorem}\label{thm:finite-no-additive}
四辺の点グラフが有限な標準有限Frobenius窓の射影安定群は加法群を含まない。
従ってその恒等成分は自明または一次元トーラスである。後者の添字保持類は
定理\ref{thm:simple-weight-finite-windows}と\ref{thm:double-middle-finite-windows}の
一曲面と二曲線で尽き、正則性・周期・重複同型も両定理の通りである。
\end{theorem}
\begin{proof}
加法群があれば定理\ref{thm:finite-connected-regularity}により既に全次数正則である。
型 $(2,1,1)$ の正準点対の六方程式を一次化する。二つの負重みの鎖の始元からは高々二本、
その $D$ 像からは零本である。実際 $xU,Ux,xx$ は両点で一次部を持たない。
重み0以上の不変二元は、重み0では $UU\oplus k(xz-zx)$ に入り、
それより高い項にも $z$ がないので、この二元の一次部は合わせて高々一本である。
六次元の周囲でJacobi階数が3以下となり、グラフの射影は閉埋込みなので接次元はちょうど3である。

型 $(3,1)$ では正準点で非零の文字 $z$、同じ鎖の $v=Nz,x=Nv$、自明文字 $t$ を使う。
二つの鎖始元 $r_1,r_2$ の一次部は高々二本である。重み $-2$ の初期成分は
$zv-vz,zt,tz$ の線型結合で、$D$ を掛けた重み0の部分は $zx-xz,vt,tv$ の線型結合となる。
よって $Dr_1,Dr_2$ の一次部は合わせて高々一本である。重み0より大きい項は $z$ を含めない。
$D^2 r_1,D^2 r_2$ はそのような項だけである。ここでも接次元3となる。
両型とも前補題に反する。

型 $(4)$ では同定理の対角還元により六関係の初期空間は通常の交代二形式全体である。
従って文字を重み $0,1,2,3$ の順に取ると、各交換子 $ij-ji$ ($i<j$) に
重みが $i+j$ より真に大きい項を足した六式が辺の基底になる。
両側で文字 $2,3$ を零にする。残る二直線の積に制限されるのは $01-10$ の持上げだけである。
他の交換子は初期部が消え、補正の重みは2より大きい。従ってこの積に曲線が残って有限性に反する。
最後に正次元安定群のトーラスの場合は前定理を適用する。
\end{proof}

\begin{proposition}\label{prop:finite-polynomial-limit-flag}
対称性を仮定せず、有限点グラフの窓が一径数群によって多項式環の四テンソルに還元されるとする。
窓と極限が異なる軌道なら、四位置を交代体積で同一視した共通の矢重みは
スカラーと正倍率を除き $0<d<e<f$（四重みは $0,d,e,f$）であり、
\[
                  e<2d,\qquad d+e<f<2e
\]
を満たす。特に重複した重みや等差四重みでは、この閉軌道の有限グラフを得ない。
\end{proposition}
\begin{proof}
極限を固定する四位置の射影作用は同じである。実際各辺は二位置の対角グラフである。
その共通の重み基底で、関係の初期空間は交換子全体となる。
各関係には対応する交換子よりも真に大きい重みの補正だけがある。
最低重みが二重以上なら、その中の対応する二文字だけに制限すると一本の交換子しか残らない。
最低を0とし重みを $0<d\le e\le f$ とする。$e\ge2d$ なら最初の二文字の積に同じ理由で
高々一式、$f\ge2e$ なら最初の三文字の積に高々三式しかなく、いずれも有限性に反する。

$f\le d+e$ と仮定する。上三文字どうしの交換子の補正は全て $d+e$ より上なので、
最低文字を含まない。この三式は最低点対で一次部を持たず、残る三式の一次部は
相手の上三文字に関して三角可逆である。よってそこは三次元接点であり、
上三文字とその三式は定理\ref{thm:maximal-tangent-ore}の内在的な三生成基底である。
良い還元によって元の窓は標準正則であるから、定理\ref{thm:maximal-tangent-complete}により
この基底の各点曲線は滑らか又は節点のみの被約平面三次曲線となる。
ところがその左評価の三行三列行列式 $F$ には重み $d+e+f$ 以下の項がない。
実際各行の最低重みは当該二文字の和で、行列式では列文字の総重み $d+e+f$ を引く。
等号の項は通常の交換子三式の行列式で、恒等零である。
上三文字を $u,v,y$ とすれば、$u^3,u^2v,u^2y,uv^2,uvy$ の重みはすべて $d+e+f$ 以下である。
従って $F$ が恒等零なら平面全体が残り、零でなければ $[1:0:0]$ は
二次接錐が零または $y^2$ の倍数となる特異点である。いずれも許容曲線でない。
以上で $f>d+e$、既得の $f<2e$ と合わせて $e>d$ も従う。
\end{proof}
